% ============================================================
%  chapters/01_introduction.tex
% ============================================================

\section{Introduction}
\label{sec:introduction}

Artificial intelligence (AI) and deep learning are now woven throughout the computational pathology workflow, supporting diagnosis, prognosis, biomarker prediction, and treatment-response assessment from whole-slide images (WSIs)~\citep{tiwari2025current, chang2026applications}. The emergence of self-supervised pathology foundation models, including UNI, CTransPath, Virchow, and Prov-GigaPath, has accelerated this adoption by providing reusable representations for dozens of downstream clinical tasks with minimal task-specific annotation~\citep{ochi2024pathology}. Throughout this review, we use ``WSI'' (singular) and ``WSIs'' (plural) as standard abbreviations for whole-slide image(s).

Yet a growing body of evidence reveals substantial performance disparities across patient subgroups. Models for breast, lung, and brain cancers exhibit 3--16\% area under the receiver operating characteristic curve (AUROC) gaps across demographic groups~\citep{vaidya2024demographic, lin2025contrastive}. Deep networks recover acquisition-site identity from TCGA slide features with up to 86\% accuracy using standard architectures~\citep{dehkharghanian2023biased}, and site-classification AUROC $>$0.96 across multiple cancer types~\citep{howard2021site, komen2024batch}. Note that accuracy and AUROC measure different aspects: accuracy reflects correct classification rate, while AUROC $>$0.96 indicates near-perfect ranking discrimination between sites. When preserved-site cross-validation is applied, the AUROC for ancestry prediction in TCGA-BRCA collapses from 0.798 to 0.507~\citep{howard2021site}. Howard et al. further reported that the majority of previously reported ancestry signals were confounded with tissue-source-site, underscoring the critical importance of preserved-site evaluation protocols. These disparities propagate into differential misdiagnosis, treatment-selection errors, and unequal patient outcomes~\citep{chen2023algorithmic, celi2022sources, norori2021addressing}, making fairness a prerequisite for clinical deployment under emerging regulatory frameworks such as the US FDA Good Machine Learning Practice principles and the EU AI Act~\citep{cheng2021challenges}. Throughout this review, we use ``TSS'' (tissue-source-site) to refer to the combined biological, technical, and institutional signatures that identify the originating laboratory or hospital.

Fairness challenges in histopathology differ fundamentally from those in other medical-imaging modalities. First, WSIs carry tissue-level confounders: fixation, staining, sectioning, and scanner hardware, which imprint measurable tissue-source-site (TSS) signatures on learned features~\citep{howard2021site, kheiri2025investigation, bidgoli2022bias}. These batch effects persist even in foundation-model representations: when features are kept frozen (not fine-tuned), linear classifiers predict hospital-of-origin with $>$90\% accuracy on TCGA-LUSC and CAMELYON16~\citep{komen2024batch}. Second, demographic bias in histopathology carries a biological component largely absent from radiology or retinal imaging; deep networks recover self-reported race from skin-histology patches with an AUROC of approximately 0.70, driven primarily by epidermal features~\citep{chen2025predicting}. Third, public benchmarks, The Cancer Genome Atlas (TCGA) above all others, dominate training corpora, coupling biological, technical, and demographic signals in ways that are difficult to disentangle~\citep{howard2021site, tasci2022bias, montezuma2025unbiased}. Taken together, these factors render modality-agnostic fairness recommendations insufficient for guiding equitable development of histopathology AI, however valuable they may be in other contexts~\citep{riccilara2022addressing, chen2023algorithmic, yang2024survey, chinta2025aidriven, musa2026cross}.

Existing reviews of fairness in medical AI draw mainly on chest radiology, dermoscopy, and retinal imaging, treating fairness as a modality-agnostic problem~\citep{riccilara2022addressing, chen2023algorithmic, celi2022sources, yang2024survey}. None have focused specifically on histopathology, nor measured how much empirical practice lags behind the recommendations these reviews advance. To close this gap, we conducted a PRISMA-2020 systematic review of 78 histopathology AI fairness studies published between 2017 and 2026. We pursue four objectives: (i)~catalog the bias categories observed in histopathology AI; (ii)~synthesize the empirical evidence on debiasing methods, including stain normalization, adversarial debiasing, fairness-aware representation learning, and foundation-model adaptation; (iii)~compare experimental findings with review-paper consensus to identify divergences; and (iv)~formulate a concrete research and reporting agenda for equitable pathology AI, including a histopathology-specific fairness benchmark proposal.

We make three contributions. First, we present the first histopathology-specific systematic review to quantify the gap between consensus recommendations and empirical practice~\citep{riccilara2022addressing, chen2023algorithmic, montezuma2025unbiased}. Second, we derive an eight-category bias taxonomy from experimental evidence, extending beyond the three- or four-category schemes common in modality-agnostic reviews~\citep{tasci2022bias, yang2024survey, howard2021site, dehkharghanian2023biased, vaidya2024demographic}. Third, we propose a draft reporting checklist and benchmark specification tailored to whole-slide image (WSI) multiple-instance learning (MIL) models~\citep{zong2023medfair, koh2021wilds, komen2024batch, huang2025knowledge}. Full inclusion and exclusion criteria, search strings, and the PRISMA flow are provided in Section~\ref{sec:methods}.
